% titel: Blatt 1, A10: Gemischtypen / A11: Trennmethoden
% skript: Mortimer Kap. 1 (Gemische, Trennverfahren)

% Teil I: Werkzeuge, kein Lösungsweg.

Stoffgemische klassifizieren (heterogen/homogen, fest/flüssig/gasförmig) und
geeignete Trennmethoden wählen.

\begin{werkzeugkasten}[Gemischtypen]
\begin{itemize}
  \item \textbf{Homogen} (eine Phase): Lösung, Nebel?, \dots{}  --- überall gleiche
        Zusammensetzung.
  \item \textbf{Heterogen} (mehrere Phasen): Gemenge, Suspension (fest in
        Flüssigkeit), Emulsion (flüssig in flüssig), Rauch (fest in Gas),
        Nebel (flüssig in Gas), Wolke, Schaum.
  \item \textbf{Kolloid vs.\ grobdispersion:} Teilchengröße entscheidet
        (Lösung = molekulare Verteilung; Suspension/Emulsion = sichtbare
        Partikel in einer anderen Phase).
\end{itemize}
\textit{Skript:} Mortimer Kap. 1.3 (Gemische).
\end{werkzeugkasten}

\begin{werkzeugkasten}[Trennmethoden]
\begin{itemize}
  \item \textbf{Filtration:} unlöslichen festen Bestandteil aus einer
        Flüssigkeit trennen (z. B.\ AgCl-Niederschlag).
  \item \textbf{Destillation:} Flüssigkeiten über Siedepunktsdifferenz trennen.
  \item \textbf{Extraktion:} gelösten Stoff in ein anderes Lösungsmittel
        überführen (Verteilung zwischen zwei Phasen).
  \item \textbf{Zentrifugation:} Dichteunterschiede über Beschleunigung
        trennen (z. B.\ Blutplasma/Zellen).
  \item \textbf{Chromatographie:} verschiedene Affinitäten zu stationärer
        Phase trennen (z. B.\ Farbstoffe).
\end{itemize}
\textit{Skript:} Mortimer Kap. 1.4 (Trennverfahren).
\end{werkzeugkasten}

\begin{achtung}[Darauf achten]
\begin{itemize}
  \item Mehrere Methoden können gleichzeitig passen --- mehr als eine Option
        ist erlaubt.
  \item Die \emph{Phase} und die \emph{Löslichkeit} des Stoffes bestimmen,
        welche Methode greift.
  \item Flüssiges Beton ist eine Suspension (Feststoff in Wasser), kein
        homogenes Gemisch.
\end{itemize}
\end{achtung}
